\documentclass[10pt,a4paper]{amsart}
\usepackage[margin=19mm]{geometry}
\usepackage{amsmath,amssymb,mathtools}
\usepackage[scr=rsfs,cal=euler]{mathalfa}
\usepackage{xfrac}
\usepackage[unicode,hidelinks,bookmarks=false]{hyperref}
\newcommand{\ie}{i.e.\ }
\newcommand{\cf}{cf.\ }
\newcommand{\resp}{respectively\ }
\newcommand{\Subsection}{\subsection}
\providecommand{\nobreakdash}{\nobreak\hbox{-}}
\setlength{\emergencystretch}{2em}
\multlinegap0pt
\usepackage{amssymb}
\newtheorem{proposition}{Proposition}
\newtheorem{theorem}{Theorem}
\title{Analytic continuation of weighted $H$-harmonic Bergman spaces}
\author{Mat\v{e}j Morav\'ik}
\date{Source: arXiv:2606.19810v1 (CC BY 4.0)}
\begin{document}
\maketitle

\noindent\emph{Source and licence.} Analytic continuation of weighted $H$-harmonic Bergman spaces. Version arXiv:2606.19810v1 (CC BY 4.0), distributed by arXiv under the Creative Commons Attribution 4.0 International licence, http://creativecommons.org/licenses/by/4.0/, as recorded in the arXiv licence metadata for this exact version and shown on the abstract page https://arxiv.org/abs/2606.19810v1. Full licence notice: \texttt{licenses/arxiv-2606.19810-CC-BY-4.0.txt}.\par\smallskip

\noindent\textit{Editorial scope.} Counted unit: the article's theorem theorem (source0.tex lines 978--1012). The article's complete original proof of it is delivered with it (source0.tex lines 1013--1107, one \texttt{proof} environment of 95 lines). No mathematics is written, repaired or re-derived: the statement and the proof are the article's own lines, in the article's own notation, and no retained line is substituted or re-flowed.  Retained uncounted as the source-local context the proof needs: lines 301--304; lines 371--405.  Nothing was omitted from any retained span, so the package carries no omission record.  No retained line contains a citation, so the wrapper carries no bibliography and no bibliography entry.  No retained line contains a figure environment, an image-inclusion command or drawing code, so no image is delivered and the package declares no asset dependency.  Editorial formatting only, in addition: an AMS \texttt{amsart} preamble that restates the article's own declarations and macros used by the retained lines, namely the environments \texttt{proposition}, \texttt{theorem} on separate counters, together with the standard packages those lines need.  The retained environments print in this document as Proposition 1, Theorem 1 in order of appearance, whereas the article prints the same items with the numbering of its own full document.  The complete native source of the article as distributed in the arXiv e-print for this exact version is bundled unchanged as \texttt{sources/203174/source0.tex}.
\begin{align}\label{TheIdentity}
   &\lim_{\epsilon \to 0} \epsilon \, {}_p F_q\left( \begin{matrix} a_1, \dots, a_p \\[3pt] b_1, \dots, b_{q-1}, -m+\epsilon \end{matrix} \;;\; z \right) \nonumber \\
   &\quad = \frac{(-1)^m z^{m+1}}{m! (m+1)!} \frac{\prod_{i=1}^p (a_i)_{m+1}}{\prod_{j=1}^{q-1} (b_j)_{m+1}} \, {}_p F_q\left( \begin{matrix} a_1+m+1, \dots, a_p+m+1 \\ b_1+m+1, \dots, b_{q-1}+m+1, m+2 \end{matrix} \;;\; z \right),
\end{align}

  \begin{proposition}\label{Pro 3.1}
      When $n$ is even we have:
      \begin{align} \label{sude}
I_m(s) &= (n-1)_m^2  \Gamma\!\left(
\begin{matrix}
\frac{n}{2},\; s+2n-1,\; m+\frac{n}{2}-1,\; s+n
\\
m,\; m+n-1,\; s+\frac{3n}{2},\; m+s+\frac{3n}{2}-1
\end{matrix}
\right) \nonumber \\[3. pt]
&\quad \times {}_4F_3 \left( \begin{matrix} \frac{n}{2}, \, s - m + \frac{n}{2} + 1, \, 1 - \frac{n}{2}, \, s + n \\ s + \frac{3n}{2}, \, 2 - m - \frac{n}{2}, \, s + \frac{n}{2} + 1 \end{matrix} ; 1 \right),
      \end{align}
      for $m >0.$\\
      When $n$ is odd we have 
\begin{align} 
I_m(s) &= \frac{(n-1)_m^2}{(n+s)_m}\,
\Gamma\!\left(
\begin{matrix}
s+\frac{n}{2}+1,\; s+2n-1,\; 1-m-\frac{n}{2}
\\
s-m+\frac{n}{2}+1,\; 1-\frac{n}{2},\; m+2n+s-1
\end{matrix}
\right) \label{liche} \\[3. pt]
&\qquad \quad \times {}_4F_3 \left( \begin{matrix} n + s, \, m + n - 1, \, m + s + \frac{3n}{2} - 1, \, m \\ m + n + s, \, m + \frac{n}{2}, \, m + 2n + s - 1 \end{matrix} ; 1 \right) \nonumber \\[1em]
& + (n-1)_m^2  \Gamma\!\left(
\begin{matrix}
\frac{n}{2},\; s+2n-1,\; m+\frac{n}{2}-1,\; s+n
\\
m,\; m+n-1,\; s+\frac{3n}{2},\; m+s+\frac{3n}{2}-1
\end{matrix}
\right) \nonumber \\[3. pt]
&\qquad \quad \times {}_4F_3 \left( \begin{matrix} \frac{n}{2}, \, s - m + \frac{n}{2} + 1, \, 1 - \frac{n}{2}, \, s + n \\ s + \frac{3n}{2}, \, 2 - m - \frac{n}{2}, \, s + \frac{n}{2} + 1 \end{matrix} ; 1 \right) \nonumber,
\end{align}
for $m >0.$
  \end{proposition}

\begin{theorem}
    For $n$ odd, we have
    \begin{align*}
  \lim_{s\rightarrow -n+\frac{1}{2}} (s+n-\frac{1}{2})\,I_m(s) &=  - \Gamma\!\left(
\begin{matrix}
m+n-1,\; m-\frac12,\; \frac n2,\; \frac n2,\; n-\frac12,\; n-\frac12
\\[3pt]
\frac12, \; m, \; m, \; \frac{n-1}{2}, \; n, \; n-1, \; n-1, \; \frac{n+1}{2}
\end{matrix}
\right)\\[3pt]
& \quad \times {}_4F_3\left( \begin{matrix} \frac{1}{2}, & \frac{n}{2}, & n-\frac{1}{2}, & 1-m \\ \frac{n+1}{2}, & n, & \frac{3}{2}-m \end{matrix} \;;\; 1 \right),
\end{align*}
and
\begin{align*}
     \lim_{s\rightarrow -n-\frac{1}{2}} (s+n+\frac{1}{2})\,I_m(s) &=
-\Gamma\!\left(
\begin{matrix}
\frac32, \; m+n-1,\; \frac n2,\; n-\frac32,\; n+\frac12,\;
m+\frac n2+\frac12,\; m-\frac32
\\[3pt]
n-1,\; n-1,\; m,\; m,\; m+\frac n2-\frac32,\;
\frac{n+1}{2},\; \frac{n+3}{2},\; 1-\frac n2,\; n
\end{matrix}
\right)
\\[3pt]
&\quad \times
{}_4F_3 \left(
\begin{matrix}
\frac32,\; \frac n2,\; n+\frac12,\; 1-m\\[3pt]
\frac{n+3}{2},\; \frac52-m,\; n
\end{matrix};1
\right).
 \end{align*}
In particular, we have  $-n-\frac{1}{2}, \, -n+\frac{1}{2} \in \mathcal{W}^d_1$.
\end{theorem}

\begin{proof}
        When $n$ is odd, from Proposition~\ref{Pro 3.1} we have
\begin{align} 
I_m(s) &= \frac{(n-1)_m^2}{(n+s)_m}\,
\Gamma\!\left(
\begin{matrix}
s+\frac{n}{2}+1,\; s+2n-1,\; 1-m-\frac{n}{2}
\\
s-m+\frac{n}{2}+1,\; 1-\frac{n}{2},\; m+2n+s-1
\end{matrix}
\right) \label{liche_thm} \\[3pt]
&\qquad \quad \times {}_4F_3 \left( \begin{matrix} n + s, \, m + n - 1, \, m + s + \frac{3n}{2} - 1, \, m \\ m + n + s, \, m + \frac{n}{2}, \, m + 2n + s - 1 \end{matrix} ; 1 \right) \nonumber \\[1em]
& + (n-1)_m^2  \Gamma\!\left(
\begin{matrix}
\frac{n}{2},\; s+2n-1,\; m+\frac{n}{2}-1,\; s+n
\\
m,\; m+n-1,\; s+\frac{3n}{2},\; m+s+\frac{3n}{2}-1
\end{matrix}
\right) \label{sude_thm} \\[3pt]
&\qquad \quad \times {}_4F_3 \left( \begin{matrix} \frac{n}{2}, \, s - m + \frac{n}{2} + 1, \, 1 - \frac{n}{2}, \, s + n \\ s + \frac{3n}{2}, \, 2 - m - \frac{n}{2}, \, s + \frac{n}{2} + 1 \end{matrix} ; 1 \right) \nonumber.
\end{align}
We first compute residue at $s = -n+\frac1 2.$
For $n > 3$ each Gamma function in \eqref{liche_thm} is regular at $s = -n + \frac{1}{2}$. For $n = 3$ we have $$\lim_{s \rightarrow-n+\frac{1}{2}}\frac{\Gamma(s+\frac{n}{2}+1)}{\Gamma(s-m+\frac{n}{2}+1)} = (-1)^m\,m!,$$
hence the first prefactor is regular for each $n \geq 3$.
Also the prefactor of \eqref{sude_thm} is regular, thus 
\begin{align*}
   \lim_{s\rightarrow -n+\frac{1}{2}} (s+n-\frac{1}{2})\,I_m(s) &=  G_m \lim_{s\rightarrow -n+\frac{1}{2}} (s+n-\frac{1}{2})\,{}_4F_3 \left( \begin{matrix} \frac{n}{2}, \, s - m + \frac{n}{2} + 1, \, 1 - \frac{n}{2}, \, s + n \\ s + \frac{3n}{2}, \, 2 - m - \frac{n}{2}, \, s + \frac{n}{2} + 1 \end{matrix} ; 1 \right)
\end{align*}
where 
$$
G_m = (n-1)_m^2\,  \Gamma\!\left(
\begin{matrix}
\frac{n}{2},\; n-\frac{1}{2},\; m+\frac{n}{2}-1,\; \frac{1}{2}
\\
m,\; m+n-1,\; \frac{n+1}{2},\; m+\frac{n}{2}-\frac{1}{2}
\end{matrix}
\right).
$$
Using \eqref{TheIdentity} we get
\begin{align}
    \lim_{s\rightarrow -n+\frac{1}{2}} (s+n-\frac{1}{2})\,&{}_4F_3 \left( \begin{matrix} \frac{n}{2}, \, s - m + \frac{n}{2} + 1, \, 1 - \frac{n}{2}, \, s + n \\ s + \frac{3n}{2}, \, 2 - m - \frac{n}{2}, \, s + \frac{n}{2} + 1 \end{matrix} ; 1 \right)\nonumber\\
    &= 
    -\frac{\Gamma(m + \frac{n-1}{2})
\Gamma\left(n-\frac12\right)
}{\Gamma(m)\,
\Gamma\left(\frac{n-1}{2}\right)
\Gamma\left(1-\frac n2\right)
\Gamma(n)
}\frac{\Gamma\!\left(2-m-\frac{n}{2}\right)}{\Gamma\!\left(\frac{3}{2}-m\right)}\\
& \times
_4F_3\left( \begin{matrix} \frac{1}{2}, & \frac{n}{2}, & n-\frac{1}{2}, & 1-m \\ \frac{n+1}{2}, & n, & \frac{3}{2}-m \end{matrix} \;;\; 1 \right),
\end{align}
here we rewrote Pochhammer symbols as Gamma functions.
Using Legendre's formula \begin{equation}\label{lagrange}
\Gamma(z)\Gamma(1-z)=\frac{\pi}{\sin(\pi z)},
\end{equation} we get 
$$
\frac{\Gamma\!\left(2-m-\frac{n}{2}\right)}{\Gamma\!\left(\frac{3}{2}-m\right)} = \frac{\Gamma\!\left(m - \frac{1}{2}\right)}{\Gamma\!\left(m + \frac{n}{2} - 1\right)}.
$$
Putting all of this together we get
\begin{align*}
  \lim_{s\rightarrow n+\frac{1}{2}} (s+n-\frac{1}{2})\,I_m(s) &=  - \Gamma\!\left(
\begin{matrix}
m+n-1,\; m-\frac12,\; \frac n2,\; \frac n2,\; n-\frac12,\; n-\frac12
\\[3pt]
\frac12, \; m, \; m, \; \frac{n-1}{2}, \; n, \; n-1, \; n-1, \; \frac{n+1}{2}
\end{matrix}
\right)\\
& \quad \times {}_4F_3\left( \begin{matrix} \frac{1}{2}, & \frac{n}{2}, & n-\frac{1}{2}, & 1-m \\ \frac{n+1}{2}, & n, & \frac{3}{2}-m \end{matrix} \;;\; 1 \right),
\end{align*}
 now as $\frac{(1-m)_{\ell}}{(\frac{3}{2}-m)_{\ell}}$ is positive for $0 \leq \ell \leq m-1$ we see that $-n+\frac{1}{2} \in\mathcal{W}^d_1$.
 
Exactly the same argument shows, that 
\begin{align*}
     \lim_{s\rightarrow -n-\frac{1}{2}} (s+n+\frac{1}{2})\,I_m(s) &=
-\Gamma\!\left(
\begin{matrix}
\frac32, \; m+n-1,\; \frac n2,\; n-\frac32,\; n+\frac12,\;
m+\frac n2+\frac12,\; m-\frac32
\\[3pt]
n-1,\; n-1,\; m,\; m,\; m+\frac n2-\frac32,\;
\frac{n+1}{2},\; \frac{n+3}{2},\; 1-\frac n2,\; n
\end{matrix}
\right)
\\[6pt]
&\quad \times
{}_4F_3 \left(
\begin{matrix}
\frac32,\; \frac n2,\; n+\frac12,\; 1-m\\[3pt]
\frac{n+3}{2},\; \frac52-m,\; n
\end{matrix};1
\right),
 \end{align*}
thus $-n-\frac{1}{2} \in \mathcal{W}_1^d$ and our conclusion follows.
\end{proof}

\end{document}
